\section{怎样分析句子成分之间的关系}

一九八一年全国高考语文试卷中有这样一道题：

“在这一斗争中，凡是拥有某种尽管是微不足道的但是有利于生存斗争的个别特质的个体，都最有希望达到成熟和繁殖。”要求回答问题：“拥有”的宾语是什么？

不少考生答道：“拥有”的宾语是“个体”。这个答案当然错了。为什么说它错了呢？这得从如何分析句子成分之间的关系谈起。

\textbf{1. 根据语序和句意，找准主语和谓语。}

这是分析句子成分之间的关系的关键，因为主语或谓语找错了，其他句子成分之间是如何互相搭配的就全部乱了套。我们在上一节中已经说过，句子成分在句子里是有一定位置的。一般说来，一个单句的语序是：

\hspace{4em}主语部分\hspace{6em}谓语部分

定语 + （的） + 主语 $\Vert$ + 状语 +（地）+ 谓语 +（得）+
补语 + 定语 +（的）+ 宾语

当然，有的状语可以放在句首（如上述高考题中的“在这一斗争中”），而且在绝大多数情况下，一个句子不会出现所有的句子成分，不过，主语和谓语则通常是少不了的。我们根据语序和句意，可以提问：全句被陈述的对象是“谁”或“什么”？这个被陈述的对象“怎么样”或“做什么”，“是什么”？例如：上述高考题中能够作为全句被陈述对象的，当然不是“在这一斗争中”，因为这是个介词结构，没有资格作主语，在这儿只能是提前的状语；也不可能是“个别特质”或“拥有······个别特质”，因为在“特质”之后还有“的个体”三个字，而“的”字是定语结构助词，这就决定了它们都没有资格作这个句子的主语；剩下的只有“个体”这个名词或“拥有……个体”这个动宾词组可能作这个句子的主语了。但是，“拥有……个体”如果作全句的主语的话，从整个句意一分析，“\uuline{拥有······个体}$\Vert$[都最]\uline{有} \uwave{希望} \uline{达到} \uwave{成熟和繁殖}”简直不知所云，主语和谓语之间不能构成“谁（什么）$\Vert$ 怎么样（做什么、是什么）”的意义关系，因此也没有资格作这个句子的主语，我们从而能够肯定只有“个体”才有资格作这个句子的主语。“个体”怎么样？“都最有希望达到成熟和繁殖。”这样一分析，全句的主语部分和谓语部分就明确起来了。

\textbf{2.借助虚词，确定中心词和修饰语之间的关系。}

一个主语或宾语前面有几个定语，一般来说，最后一个“的”字前面的是定语，后面的是主语或宾语。根据这一点，也可以确定“······的个别特质的个体”中的“个别特质”只能在定语范围内，“个体”才是主语部分的中心词；如果一个谓语前面有几个状语，一般来说，最后一个状语结构助词“地”字前面的是状语，后面的动词或形容词是谓语，如“全国各族人民热烈地日思夜想地向往首都北京城。”在这儿“日思夜想”仍然是“向往”的状语；当然补语结构助词“得”字前面的是谓语，后面的是补语。

介词和宾语结合所构成的介词结构不能充当主语、谓语或宾语，只能充当定语、状语或补语。副词也只能充当状语或补语。关联词语只能起连结词语、词组或者分句的作用，而所连接的词语、词组或者分句在全句中作什么成分，得出它们的位置和意义来决定，例如“尽管是微不足道的，但是有利于生存斗争的”，孤立地看它是由转折关联词语连接的两个分句，但是在上述高考试题中，它只能作“个别特质”的定语成分。

我们只要恰当地借助这些虚词，再分析修饰和被修饰、限制和被限制的关系，就能找出句子的主干（主语、谓语和宾语），区别出修饰或限制主干成分的附加成分。如开头所述的那道高考试题的中心词和修饰语的关系，借助虚词，可以明确区分如下：

[在这一斗争中]，凡是（
\uline{拥有}
某种尽管是微不足道的但是有利于生存斗争的个别
\uwave{特质}
）
的
\uuline{个体}
$\Vert$
都［最］
\uline{有}
\uwave{希望}
\uline{达到}
\uwave{成熟}
和
\uwave{繁殖}。

\textbf{3. 运用代换法逐层剖析，弄清各个成分之间的关系。}

找出了句子的主干成分，明确了句子的基本结构之后，就可以逐一确定附加成分。由单词作附加成分的比较好办，如上述那道试题中的“最有希望”的“最”，很容易断定是作修饰动词谓语的“有”的状语用的；“凡是”、“都”则要复杂一点，不容易一眼看出，其实这是两个关联词，不作句子成分，在这儿它们只起强调的作用。我们可以用“代换法”提问：“凡是”什么？答案只能是“个体”。它与“都”呼应，强调“个体”的条件，而不作“拥有”的状语。问题就难在用词组甚至复句结构作附加成分上，这也正是一个单句复杂之所在。比如上述那道试题中的主语“个体”的附加成分是：“凡是拥有某种尽管是微不足道的但是有利于生存斗争的个别特质”，这是一个比较复杂的定语。“凡是”是强调“个体”的条件关系的关联词，可以不再讨论外，我们借助这个复杂定语中最后的一个结构助词“的”，确定“（个别）特质”的定语是“某种”和“尽管······但是······”这个复句结构，而整个定语的前面又有动词“拥有”，就可断定“（个别）特质”是“拥有”的宾语；或者用“代换法”提问：“拥有”什么？很明显，就是拥有“特质”，因为若是认为“拥有”“个体”的话，整个句意就无法理解了。这个复杂定语的成分关系是：

\uline{拥有}（某种）（尽管是微不足道的但是有利于生存斗争的）（个别）\uwave{特质}······

说到这里，我们就可以明白为什么说把“个体”看作“拥有”的宾语是错的了。

\textbf{4. 根据各种句式的特点，确定句子成分之间的关系。}
汉语句式多种多样，常见的有省略式、主谓倒置式，以及兼语式、连动式等等，我们必须根据语序和句意仔细地加以辨析，才能确定句子成分之间的关系。例如：

①\ 仿佛一张极大极大的荷叶铺着，满是奇异的绿呀。

（“铺着”的主语是什么？应该是“绿”，因为这是一个主谓倒置句。）

②\ 我也想警告某些人，当心呻吟着的那些锭子上的冤魂！

（“当心”的主语是什么？应该是“某些人”，因为这是一个兼语式句子；“某些人”是“警告”的宾语，又是“当心”的主语。

“呻吟着”修饰的中心词是什么？应该是“冤魂”，而不是“那些锭子上”；“呻吟着”与“那些锭子上”同是“冤魂”的修饰语，否则语意不通。）

③\ 你会看见沿途汇入千百条泉流，逐渐形成溪流，再汇入许多涧流和溪流，就形成河流，奔腾出天山。

（“看见”的宾语是什么？应该是“看见”以下的全部文字，而不是其中的任何一句话。这是由许多动宾词组构成一个复杂词组作宾语的一个例子。

“形成”的主语是什么？应该是“泉流”，它承前省略了。

“出”的主语是什么？应该是“河流”，也是承前省略了。）

④\ 今晚我在院子里坐着乘凉，忽然想起日日走过的荷塘，在这满月的光里，总该另有一番样子吧。

（“我”的谓语是什么？应该是“坐着”“乘凉”和“想起”，“想起”的主语是承前省略了。与“荷塘”搭配的谓语是什么？应该是“想起”，而不是“走过”，“走过”在这儿是“荷塘”的定语。“有”的主语是什么？应该是“荷塘”，而不能是“在这满月的光里”，因为介词结构不能作主语。）

⑤\ 我曾以博物学者的资格参加贝格尔号巡洋舰的环球远航，在南美洲看到的关于生物的地理分布和现存生物与古生物在地质上的关系，给了我很深刻的印象。

（介词“以”的宾语是什么？应该是“资格”，而不是“博物学者”，“博物学者”是“资格”的定语。“在南美洲看到”这个偏正词组所修饰的中心词是什么？应该是“关系”。“关于生物……在地质上”也是修饰“关系”的定语；不能认为“在南美洲看到”所修饰的中心词是“地理分布”，或者是“地理分布”和“关系”。）

（“给”的主语是什么？应该是“关系”，而不是别的什么词语。）

⑥\ 本身具备着可能发展条件的人类的远祖，正是在一定的环境条件下从古猿分化出来之后，通过必需的生活活动，使前肢解放为手，用双手制造并使用工具来改造自然，在改造自然的进程中逐步改造了自身，终于由接近类人猿的原始人发展成为现代人。

（“具备着”的宾语是什么？应该是“条件”，而非“人类”，当然更不是“远祖”。

“使”的状语是什么？应该是“在······下”、“从······之后”、“通过······活动”三个介词结构，而不只是其中的任何一个。

“用双手制造并使用工具来改造自然”所省略的主语是什么？应该是“远祖”，承前省略了。这是一个由兼语式和连动式构成的一个省略句，可以用划线法分析其句子成分如下：
\uline{用}
双手
\uline{制造}
并
\uline{使用}
\uwave{工具}
\uline{来改造}
\uwave{自然}。）

介词“由”的宾语是什么？应该是“原始人”，而不是“类人猿”，更不是“接近类人猿”）。

⑦\ 真正刻苦修养，忠实作马克思列宁主义创始人的学生的人，他们特别注重的，是要象马克思列宁主义创始人那样，站在马克思列宁主义立场，用马克思列宁主义的观点和方法，去解决无产阶级领导的革命运动中的各种问题。

（整个句子的主语是什么？是“他们特别注意的”，而不是“他们”，“他们”只与“特别注意”构成主谓词组，作整个“的字结构”的组成部分。

“作”的宾语是什么？应该是“学生”，而不是“人”，更不是“马克思列宁主义创始人”。 “无产阶级领导”所要修饰的中心词是什么？应该是“革命运动”，而不是“问题”。）

⑧\ 美国白皮书和艾奇逊信件的发表是值得庆祝的，因为它给了中国怀有旧民主主义思想亦即民主个人主义思想，而对人民民主主义，即民主集体主义，或民主集中主义，或集体英雄主义，或国际主义的爱国主义，不赞成，或不甚赞成，不满或有些不满，甚至抱有反感，但是还有爱国心，并非国民党反动派的人们，浇了一瓢冷水，丢了他们的脸。

（“它”代的是什么？应该是指“美国白皮书和艾奇逊信件的发表”，而不是“美国白皮书”或者“美国白皮书和艾奇逊信件”。

“人们”的定语是什么？应该是“中国怀有旧民主主义思想……并非国民党反动派”这一大串话，而不仅是其中任何的一个。

“浇”的主语是什么？应该是“它”，而不是“人们”，也不是“美国白皮书”或者“美国白皮书和艾奇逊信件”。）

总之，我们只要能恰当运用句子成分分析法，联系上下文意，并综合运用词语搭配、句式特点等有关常识和辨析方法，完全可以准确地辨析句子成分之间的关系。
\newpage